% chktex-file 0
\documentclass{ctexart}
\usepackage{graphicx} % Required for inserting images
\usepackage[margin=2.5cm]{geometry}
\usepackage{booktabs}
\usepackage{float}
\usepackage{soulutf8}
\usepackage{enumitem}
\usepackage{pifont}
\usepackage{pgfplots}
\usepackage{longtable}
\pgfplotsset{compat=newest}
\usepackage{tikz}
\usepackage{xcolor} % 用于颜色设置
\usepackage{listings} % 用于代码块
\usepackage{fontspec} % 关键：用于使用系统字体 (XeLaTeX/LuaLaTeX)
% Linux 环境中没有 Windows 宋体时，使用可用的中文衬线字体作为宋体替代
\setCJKmainfont{Noto Serif CJK SC}
\usepackage{authblk}
\usepackage[
    colorlinks=true,      % 使用颜色而非方框
    linkcolor=blue,       % 内部链接颜色
    citecolor=green,      % 引用链接颜色
    urlcolor=red,         % URL 链接颜色
    bookmarks=true,       % 生成 PDF 书签
    bookmarksopen=true,   % 默认展开书签
    pdfstartview=FitH     % 页面宽度适应窗口
]{hyperref}

\sethlcolor{yellow}%设置所需高亮颜色
\usepackage[most]{tcolorbox}
\tcbuselibrary{listings,skins,breakable}
% 页眉页脚
\usepackage{fancyhdr}
\pagestyle{fancy}

% 定义草绿色
\definecolor{grassgreen}{RGB}{124,179,66}
% 补全缺失的浅蓝色定义（防止编译报错）
\definecolor{lightblue}{RGB}{30,144,255} 

% 清空默认
\fancyhf{}

% 页眉
\fancyhead[L]{\textcolor{lightblue}{\kaishu ICS}}
\fancyhead[R]{\textcolor{lightblue}{\textit{2026.9-2027.1}}}

% 页脚页码居中
\fancyfoot[C]{\textcolor{lightblue}{\thepage}}

% 页眉、页脚线颜色
\renewcommand{\headrulewidth}{1pt}
\renewcommand{\footrulewidth}{1pt}
\renewcommand{\headrule}{\color{blue!30!white}\hrule height 1pt width\headwidth}
\renewcommand{\footrule}{\color{blue!30!white}\hrule height 1pt width\headwidth}
\newtcolorbox{fancybox}[2][blue]{ % 可选颜色参数
  enhanced,
  breakable,
    colback=#1!3!white,
    colframe=#1!18!white,
    boxrule=0.8pt,
    arc=6pt,
    outer arc=6pt,
  left=10pt,right=10pt,top=12pt,bottom=10pt,
  title={\textbf{#2}},
    fonttitle=\bfseries\color{#1!65!black},
  attach boxed title to top left={yshift=-2mm,xshift=5mm},
  boxed title style={
        colback=#1!8!white,
        colframe=#1!18!white,
        boxrule=0.6pt,
        arc=4pt,
  },
    shadow={1.5pt}{-1.5pt}{0pt}{black!10},
}

% 定义一套适合C++的配色方案
\usepackage{tabularx}
\newfontfamily{\consolasfont}{Ubuntu}

\newcommand{\bluecode}[1]{{\color{blue}\consolasfont #1}}
\lstdefinestyle{MyCppStyle}{
    language = C++,
    % 确保这里使用 \consolasfont
    basicstyle = \consolasfont, 
    keywordstyle = \color{blue!70!black}, 
    commentstyle = \color{green!50!black}, 
    stringstyle = \color{red!60!black}, 
    numberstyle = \tiny\color{gray}, 
    numbers = left, 
    frame = single, 
    rulecolor = \color{lightgray}, 
    breaklines = true, 
    backgroundcolor = \color{white}, 
    tabsize = 4, 
    showstringspaces = false, 
    xleftmargin = 2em, 
    escapeinside=``,
}
% --- lstdefinestyle HTML 风格定义 ---
\lstdefinestyle{mhtml}{
    language = HTML,
    basicstyle = \consolasfont\small, 
    tagstyle = \color{blue!70!black}\bfseries,
    keywordstyle = \color{purple!80!black},
    commentstyle = \color{green!50!black}\itshape, 
    stringstyle = \color{red!60!black}, 
    numberstyle = \tiny\color{gray}, 
    numbers = left, 
    frame = single, 
    rulecolor = \color{lightgray}, 
    breaklines = true, 
    backgroundcolor = \color{white}, 
    tabsize = 4, 
    showstringspaces = false, 
    xleftmargin = 2em, 
    escapeinside=``,
    morekeywords={als, doctype, html, head, title, meta, link, body, div, span, h1, h2, h3, h4, h5, h6, p, a, img, ul, ol, li, table, tr, th, td, pre, code, script, style},
    morekeywords=[2]{class, id, href, src, rel, type, charset, lang, name, content, style},
    keywordstyle=[2]\color{purple!80!black},
}
\usepackage{hyperref}
\hypersetup{
    colorlinks=true,      
    linkcolor=blue,       
    filecolor=magenta,    
    urlcolor=cyan,        
    pdftitle={我的文档标题}, 
}

% ==================== 正文部分 ====================
\begin{document}

% 开启无页眉页脚模式（针对封面和目录）
\pagestyle{empty}

% ------------------ 美化封面页 ------------------
\begin{titlepage}
    \centering
    % 顶部 TikZ 几何几何装饰条（全宽 bleed 效果）
    \begin{tikzpicture}[remember picture, overlay]
        \fill[blue!50!white] (current page.north west) rectangle ([yshift=-2cm]current page.north east);
        \fill[grassgreen] (current page.north west) ++(0,-2cm) rectangle ([yshift=-2.2cm]current page.north east);
    \end{tikzpicture}
    
    \vspace*{3.5cm}
    
    % 大标题
    {\fontsize{30pt}{36pt}\selectfont\bfseries\color{blue!80!black} 计算机系统导论（ICS）}\\[0.6em]
    % 副标题或英文标识
    {\fontsize{14pt}{18pt}\selectfont\color{gray}\textsf{Introduction to Computer System}}\\[1cm]
    
    % 居中装饰线
    \centering\makebox[4cm]{\color{grassgreen}\rule{5cm}{1.5pt}}
    
    \vfill
    
    % 使用 tcolorbox 包装的精致作者信息框
    \begin{tcolorbox}[
        enhanced,
        width=0.55\textwidth,
        colback=blue!5!white,
        colframe=blue!5!white,
        boxrule=1.2pt,
        arc=6pt,
        sharp corners=downhill, % 别致的对角切角效果
        shadow={2pt}{-2pt}{0pt}{black!20},
        halign=center,
        valign=center,
        top=15pt, bottom=15pt
    ]
        {\large\bfseries\color{black!80} Author：} {\large\kaishu Simson} \\[0.6em]
        {\large\bfseries\color{black!80} Date：} {\large\kaishu September 2026}
    \end{tcolorbox}
    
    \vfill
    \vspace*{2cm}
    
    % 底部标识
    {\color{gray!80}\small\the\year\ \copyright\ All Rights Reserved.}
\end{titlepage}

\pagestyle{fancy}

\newpage
\tableofcontents
\newpage

% 设置从当前页（正文第一页）开始采用阿拉伯数字编码，并自动重置为第 1 页
\pagenumbering{arabic}


\section{Lecture 1}

1. 配置Xlab

\begin{lstlisting}
    ssh-keygen -t ed25519
    result:ssh-ed25519 AAAAC3NzaC1lZDI1NTE5AAAAICu5/higBBmBl35SYAtn03EFjHaF+DA1eg7kHlxqtn7U simson@SImson-SHY-HP
\end{lstlisting}

\end{document}